% =====================================================================
%  preambulo.tex
%  Preámbulo común a TODAS las entregas (capítulos sueltos y libro final).
%  Se carga con % =====================================================================
%  preambulo.tex
%  Preámbulo común a TODAS las entregas (capítulos sueltos y libro final).
%  Se carga con % =====================================================================
%  preambulo.tex
%  Preámbulo común a TODAS las entregas (capítulos sueltos y libro final).
%  Se carga con % =====================================================================
%  preambulo.tex
%  Preámbulo común a TODAS las entregas (capítulos sueltos y libro final).
%  Se carga con \input{config/preambulo} desde la raíz del repositorio.
%
%  MOTOR: XeLaTeX (obligatorio). Es el que permite usar la fuente real
%  Times New Roman del sistema. NO compilar con pdflatex.
% =====================================================================

% ---------- Tamaño base exacto: 11 puntos ----------
% Dos precisiones importantes:
%  1) La opción [11pt] de la clase da en realidad 10.95 pt. El paquete
%     fontsize regenera todas las órdenes de tamaño con el valor exacto.
%  2) El "punto" de TeX mide 1/72.27 de pulgada, pero el punto que usan
%     Word y los visores de PDF (el que se verificará) mide 1/72 = 1 bp.
%     Por eso todos los tamaños de este proyecto se declaran en "bp":
%     así el PDF reporta 11.00, 12.00, 14.00 y 16.00 exactos.
\usepackage[fontsize=11bp]{fontsize}

% ---------- Fuentes ----------
\usepackage{amsmath, amssymb}
\usepackage{fontspec}
\defaultfontfeatures{Ligatures=TeX}

% Times New Roman del sistema (macOS y Windows la traen).
% Si no existe, se usa TeX Gyre Termes: clon libre métricamente idéntico,
% así el documento compila igual en cualquier máquina sin cambiar de aspecto.
\IfFontExistsTF{Times New Roman}{
    \setmainfont{Times New Roman}
    \newcommand{\FuenteTexto}{Times New Roman}
}{
    \setmainfont{texgyretermes}[
        Extension      = .otf,
        UprightFont    = *-regular,
        BoldFont       = *-bold,
        ItalicFont     = *-italic,
        BoldItalicFont = *-bolditalic]
    \newcommand{\FuenteTexto}{TeX Gyre Termes (clon métrico de Times New Roman)}
}
% Monoespaciada para el código
\IfFontExistsTF{Courier New}{
    % Sin Scale: así \texttt en línea mide exactamente los 11 bp del texto.
    \setmonofont{Courier New}
}{
    \setmonofont{texgyrecursor}[
        Extension      = .otf,
        UprightFont    = *-regular,
        BoldFont       = *-bold,
        ItalicFont     = *-italic,
        BoldItalicFont = *-bolditalic]
}
% Matemáticas en Times (Termes Math es la versión matemática de Times).
% Se carga por nombre de archivo: en macOS fontconfig no indexa las fuentes
% de TeX Live, así que buscarla por nombre comercial falla.
\usepackage{unicode-math}
\setmathfont{texgyretermes-math.otf}

% ---------- Idioma ----------
% es-lcroman: evita que babel componga los números romanos en versalitas.
% Times New Roman no tiene versalitas reales y XeTeX las simularía
% encogiendo la letra a ~8 pt, rompiendo el requisito de tamaños.
\usepackage[spanish, es-tabla, es-noshorthands, es-lcroman]{babel}
\usepackage{csquotes}

% ---------- Página ----------
\usepackage{geometry}
\geometry{top=2.5cm, bottom=2.5cm, left=2.5cm, right=2.5cm}
\usepackage{microtype}
\linespread{1.15}
\setlength{\parskip}{0.5em}

% ---------- Utilidades ----------
\usepackage{ifthen}
\usepackage{xcolor}
\usepackage{graphicx}
\usepackage{float}
\usepackage{booktabs}
\usepackage{longtable}
\usepackage{multirow}
\usepackage{array}
\usepackage{siunitx}
\usepackage{caption}
\usepackage{subcaption}
\usepackage{enumitem}
\usepackage{tcolorbox}
\usepackage[linesnumbered, ruled, vlined]{algorithm2e}
\usepackage[nottoc, notbib]{tocbibind}   % LOF y LOT aparecen en el índice
\usepackage{fancyhdr}
\usepackage{titlesec}
\usepackage{listings}

% Rutas donde LaTeX buscará imágenes (funciona igual compilando el libro
% completo o un capítulo suelto, porque siempre se compila desde la raíz).
\graphicspath{%
    {img/}%
    {capitulos/cap01/figuras/}%
    {capitulos/cap02/figuras/}%
    {capitulos/cap03/figuras/}%
    {capitulos/cap04/figuras/}%
    {capitulos/cap05/figuras/}%
    {capitulos/cap06/figuras/}%
    {capitulos/cap07/figuras/}%
}

% ---------- Colores institucionales ----------
\definecolor{guindaIPN}{RGB}{110, 28, 52}
\definecolor{azulESCOM}{RGB}{0, 81, 158}
\definecolor{grisCodigo}{RGB}{245, 245, 245}
\definecolor{verdeComent}{RGB}{0, 120, 60}
\definecolor{moradoKey}{RGB}{120, 20, 120}

% =====================================================================
%  JERARQUÍA TIPOGRÁFICA REQUERIDA
%     texto        11 pt
%     capítulo     16 pt negritas
%     sección      14 pt negritas
%     subsección   12 pt negritas
%     subsubsec.   12 pt negritas cursiva
%  Los tamaños se fijan de forma absoluta con \fontsize{pt}{interlínea},
%  no con \Large/\huge, para que sean exactamente los pedidos.
% =====================================================================
\newcommand{\TamCapitulo}{\fontsize{16bp}{19.2bp}\selectfont}
\newcommand{\TamSeccion}{\fontsize{14bp}{16.8bp}\selectfont}
\newcommand{\TamSubseccion}{\fontsize{12bp}{14.4bp}\selectfont}
\newcommand{\TamTexto}{\fontsize{11bp}{13.2bp}\selectfont}

\titleformat{\chapter}[hang]
    {\normalfont\bfseries\TamCapitulo}{\chaptertitlename\ \thechapter.}{0.6em}{}
\titlespacing*{\chapter}{0pt}{0pt}{24pt}

\titleformat{\section}
    {\normalfont\bfseries\TamSeccion}{\thesection}{0.6em}{}
\titlespacing*{\section}{0pt}{16pt}{8pt}

\titleformat{\subsection}
    {\normalfont\bfseries\TamSubseccion}{\thesubsection}{0.6em}{}
\titlespacing*{\subsection}{0pt}{12pt}{6pt}

\titleformat{\subsubsection}
    {\normalfont\bfseries\itshape\TamSubseccion}{\thesubsubsection}{0.6em}{}
\titlespacing*{\subsubsection}{0pt}{10pt}{5pt}

% Encabezados de los índices (son \chapter*) al mismo tamaño de 16 pt
\renewcommand{\contentsname}{Índice de contenido}

% Pies de figura y tabla en el tamaño del texto
\captionsetup{font=normalsize, labelfont=bf, labelsep=period, skip=8pt}

% ---------- Encabezados y pies ----------
\setlength{\headheight}{16pt}
\pagestyle{fancy}
\fancyhf{}
\renewcommand{\headrulewidth}{0.4pt}
\fancyhead[L]{\TamTexto\nouppercase{\leftmark}}
\fancyhead[R]{\TamTexto\Materia}
\fancyfoot[C]{\TamTexto\thepage}
\fancypagestyle{plain}{\fancyhf{}\fancyfoot[C]{\TamTexto\thepage}\renewcommand{\headrulewidth}{0pt}}

% ---------- Código fuente (C, OpenMP, MPI, CUDA) ----------
% El código va en monoespaciada a 10 pt: a 11 pt las líneas largas no caben
% en el ancho de página. El texto redactado sí respeta los 11 pt.
\lstset{
    basicstyle=\ttfamily\fontsize{10bp}{12bp}\selectfont,
    keywordstyle=\color{moradoKey}\bfseries,
    commentstyle=\color{verdeComent}\itshape,
    stringstyle=\color{azulESCOM},
    backgroundcolor=\color{grisCodigo},
    frame=single,
    rulecolor=\color{gray!50},
    numbers=left,
    numberstyle=\fontsize{8bp}{9.6bp}\selectfont\color{gray},
    numbersep=8pt,
    breaklines=true,
    breakatwhitespace=true,
    showstringspaces=false,
    tabsize=4,
    captionpos=b,
    columns=flexible,
}
% Lenguajes derivados
\lstdefinelanguage{CUDA}[ANSI]{C}{
    morekeywords={__global__, __device__, __host__, __shared__, __constant__,
                  __syncthreads, blockIdx, blockDim, threadIdx, gridDim,
                  cudaMalloc, cudaFree, cudaMemcpy, cudaDeviceSynchronize}
}
\lstdefinestyle{C}{language=C, morekeywords={\#pragma, omp, parallel, for, MPI_Init, MPI_Comm_rank, MPI_Comm_size, MPI_Send, MPI_Recv, MPI_Finalize}}
% Nombre en español para los listados de código
\renewcommand{\lstlistingname}{Código}

% ---------- Cajas de nota ----------
\newtcolorbox{nota}[1][Nota]{colback=azulESCOM!5, colframe=azulESCOM!70,
    fonttitle=\bfseries, title=#1, boxrule=0.6pt, arc=2pt}
\newtcolorbox{definicion}[1][Definición]{colback=guindaIPN!5, colframe=guindaIPN!70,
    fonttitle=\bfseries, title=#1, boxrule=0.6pt, arc=2pt}

% ---------- Bibliografía IEEE ----------
\usepackage[backend=biber, style=ieee, sorting=none, maxbibnames=6, minbibnames=3]{biblatex}
\addbibresource{bib/referencias.bib}
\DefineBibliographyStrings{spanish}{references = {Referencias}}

% ---------- Hipervínculos (cargar casi al final) ----------
\usepackage[hidelinks, bookmarksnumbered]{hyperref}
\usepackage{bookmark}

% ---------- Macros de conveniencia ----------
% Bibliografía de la entrega (aparece en el índice de contenido)
\newcommand{\Referencias}{%
    \cleardoublepage
    \printbibliography[heading=bibintoc, title={Referencias}]%
}

% Figura estándar: \Figura{archivo}{ancho}{pie de figura}{etiqueta}
\newcommand{\Figura}[4]{%
    \begin{figure}[H]
        \centering
        \includegraphics[width=#2\textwidth]{#1}
        \caption{#3}
        \label{fig:#4}
    \end{figure}%
}
 desde la raíz del repositorio.
%
%  MOTOR: XeLaTeX (obligatorio). Es el que permite usar la fuente real
%  Times New Roman del sistema. NO compilar con pdflatex.
% =====================================================================

% ---------- Tamaño base exacto: 11 puntos ----------
% Dos precisiones importantes:
%  1) La opción [11pt] de la clase da en realidad 10.95 pt. El paquete
%     fontsize regenera todas las órdenes de tamaño con el valor exacto.
%  2) El "punto" de TeX mide 1/72.27 de pulgada, pero el punto que usan
%     Word y los visores de PDF (el que se verificará) mide 1/72 = 1 bp.
%     Por eso todos los tamaños de este proyecto se declaran en "bp":
%     así el PDF reporta 11.00, 12.00, 14.00 y 16.00 exactos.
\usepackage[fontsize=11bp]{fontsize}

% ---------- Fuentes ----------
\usepackage{amsmath, amssymb}
\usepackage{fontspec}
\defaultfontfeatures{Ligatures=TeX}

% Times New Roman del sistema (macOS y Windows la traen).
% Si no existe, se usa TeX Gyre Termes: clon libre métricamente idéntico,
% así el documento compila igual en cualquier máquina sin cambiar de aspecto.
\IfFontExistsTF{Times New Roman}{
    \setmainfont{Times New Roman}
    \newcommand{\FuenteTexto}{Times New Roman}
}{
    \setmainfont{texgyretermes}[
        Extension      = .otf,
        UprightFont    = *-regular,
        BoldFont       = *-bold,
        ItalicFont     = *-italic,
        BoldItalicFont = *-bolditalic]
    \newcommand{\FuenteTexto}{TeX Gyre Termes (clon métrico de Times New Roman)}
}
% Monoespaciada para el código
\IfFontExistsTF{Courier New}{
    % Sin Scale: así \texttt en línea mide exactamente los 11 bp del texto.
    \setmonofont{Courier New}
}{
    \setmonofont{texgyrecursor}[
        Extension      = .otf,
        UprightFont    = *-regular,
        BoldFont       = *-bold,
        ItalicFont     = *-italic,
        BoldItalicFont = *-bolditalic]
}
% Matemáticas en Times (Termes Math es la versión matemática de Times).
% Se carga por nombre de archivo: en macOS fontconfig no indexa las fuentes
% de TeX Live, así que buscarla por nombre comercial falla.
\usepackage{unicode-math}
\setmathfont{texgyretermes-math.otf}

% ---------- Idioma ----------
% es-lcroman: evita que babel componga los números romanos en versalitas.
% Times New Roman no tiene versalitas reales y XeTeX las simularía
% encogiendo la letra a ~8 pt, rompiendo el requisito de tamaños.
\usepackage[spanish, es-tabla, es-noshorthands, es-lcroman]{babel}
\usepackage{csquotes}

% ---------- Página ----------
\usepackage{geometry}
\geometry{top=2.5cm, bottom=2.5cm, left=2.5cm, right=2.5cm}
\usepackage{microtype}
\linespread{1.15}
\setlength{\parskip}{0.5em}

% ---------- Utilidades ----------
\usepackage{ifthen}
\usepackage{xcolor}
\usepackage{graphicx}
\usepackage{float}
\usepackage{booktabs}
\usepackage{longtable}
\usepackage{multirow}
\usepackage{array}
\usepackage{siunitx}
\usepackage{caption}
\usepackage{subcaption}
\usepackage{enumitem}
\usepackage{tcolorbox}
\usepackage[linesnumbered, ruled, vlined]{algorithm2e}
\usepackage[nottoc, notbib]{tocbibind}   % LOF y LOT aparecen en el índice
\usepackage{fancyhdr}
\usepackage{titlesec}
\usepackage{listings}

% Rutas donde LaTeX buscará imágenes (funciona igual compilando el libro
% completo o un capítulo suelto, porque siempre se compila desde la raíz).
\graphicspath{%
    {img/}%
    {capitulos/cap01/figuras/}%
    {capitulos/cap02/figuras/}%
    {capitulos/cap03/figuras/}%
    {capitulos/cap04/figuras/}%
    {capitulos/cap05/figuras/}%
    {capitulos/cap06/figuras/}%
    {capitulos/cap07/figuras/}%
}

% ---------- Colores institucionales ----------
\definecolor{guindaIPN}{RGB}{110, 28, 52}
\definecolor{azulESCOM}{RGB}{0, 81, 158}
\definecolor{grisCodigo}{RGB}{245, 245, 245}
\definecolor{verdeComent}{RGB}{0, 120, 60}
\definecolor{moradoKey}{RGB}{120, 20, 120}

% =====================================================================
%  JERARQUÍA TIPOGRÁFICA REQUERIDA
%     texto        11 pt
%     capítulo     16 pt negritas
%     sección      14 pt negritas
%     subsección   12 pt negritas
%     subsubsec.   12 pt negritas cursiva
%  Los tamaños se fijan de forma absoluta con \fontsize{pt}{interlínea},
%  no con \Large/\huge, para que sean exactamente los pedidos.
% =====================================================================
\newcommand{\TamCapitulo}{\fontsize{16bp}{19.2bp}\selectfont}
\newcommand{\TamSeccion}{\fontsize{14bp}{16.8bp}\selectfont}
\newcommand{\TamSubseccion}{\fontsize{12bp}{14.4bp}\selectfont}
\newcommand{\TamTexto}{\fontsize{11bp}{13.2bp}\selectfont}

\titleformat{\chapter}[hang]
    {\normalfont\bfseries\TamCapitulo}{\chaptertitlename\ \thechapter.}{0.6em}{}
\titlespacing*{\chapter}{0pt}{0pt}{24pt}

\titleformat{\section}
    {\normalfont\bfseries\TamSeccion}{\thesection}{0.6em}{}
\titlespacing*{\section}{0pt}{16pt}{8pt}

\titleformat{\subsection}
    {\normalfont\bfseries\TamSubseccion}{\thesubsection}{0.6em}{}
\titlespacing*{\subsection}{0pt}{12pt}{6pt}

\titleformat{\subsubsection}
    {\normalfont\bfseries\itshape\TamSubseccion}{\thesubsubsection}{0.6em}{}
\titlespacing*{\subsubsection}{0pt}{10pt}{5pt}

% Encabezados de los índices (son \chapter*) al mismo tamaño de 16 pt
\renewcommand{\contentsname}{Índice de contenido}

% Pies de figura y tabla en el tamaño del texto
\captionsetup{font=normalsize, labelfont=bf, labelsep=period, skip=8pt}

% ---------- Encabezados y pies ----------
\setlength{\headheight}{16pt}
\pagestyle{fancy}
\fancyhf{}
\renewcommand{\headrulewidth}{0.4pt}
\fancyhead[L]{\TamTexto\nouppercase{\leftmark}}
\fancyhead[R]{\TamTexto\Materia}
\fancyfoot[C]{\TamTexto\thepage}
\fancypagestyle{plain}{\fancyhf{}\fancyfoot[C]{\TamTexto\thepage}\renewcommand{\headrulewidth}{0pt}}

% ---------- Código fuente (C, OpenMP, MPI, CUDA) ----------
% El código va en monoespaciada a 10 pt: a 11 pt las líneas largas no caben
% en el ancho de página. El texto redactado sí respeta los 11 pt.
\lstset{
    basicstyle=\ttfamily\fontsize{10bp}{12bp}\selectfont,
    keywordstyle=\color{moradoKey}\bfseries,
    commentstyle=\color{verdeComent}\itshape,
    stringstyle=\color{azulESCOM},
    backgroundcolor=\color{grisCodigo},
    frame=single,
    rulecolor=\color{gray!50},
    numbers=left,
    numberstyle=\fontsize{8bp}{9.6bp}\selectfont\color{gray},
    numbersep=8pt,
    breaklines=true,
    breakatwhitespace=true,
    showstringspaces=false,
    tabsize=4,
    captionpos=b,
    columns=flexible,
}
% Lenguajes derivados
\lstdefinelanguage{CUDA}[ANSI]{C}{
    morekeywords={__global__, __device__, __host__, __shared__, __constant__,
                  __syncthreads, blockIdx, blockDim, threadIdx, gridDim,
                  cudaMalloc, cudaFree, cudaMemcpy, cudaDeviceSynchronize}
}
\lstdefinestyle{C}{language=C, morekeywords={\#pragma, omp, parallel, for, MPI_Init, MPI_Comm_rank, MPI_Comm_size, MPI_Send, MPI_Recv, MPI_Finalize}}
% Nombre en español para los listados de código
\renewcommand{\lstlistingname}{Código}

% ---------- Cajas de nota ----------
\newtcolorbox{nota}[1][Nota]{colback=azulESCOM!5, colframe=azulESCOM!70,
    fonttitle=\bfseries, title=#1, boxrule=0.6pt, arc=2pt}
\newtcolorbox{definicion}[1][Definición]{colback=guindaIPN!5, colframe=guindaIPN!70,
    fonttitle=\bfseries, title=#1, boxrule=0.6pt, arc=2pt}

% ---------- Bibliografía IEEE ----------
\usepackage[backend=biber, style=ieee, sorting=none, maxbibnames=6, minbibnames=3]{biblatex}
\addbibresource{bib/referencias.bib}
\DefineBibliographyStrings{spanish}{references = {Referencias}}

% ---------- Hipervínculos (cargar casi al final) ----------
\usepackage[hidelinks, bookmarksnumbered]{hyperref}
\usepackage{bookmark}

% ---------- Macros de conveniencia ----------
% Bibliografía de la entrega (aparece en el índice de contenido)
\newcommand{\Referencias}{%
    \cleardoublepage
    \printbibliography[heading=bibintoc, title={Referencias}]%
}

% Figura estándar: \Figura{archivo}{ancho}{pie de figura}{etiqueta}
\newcommand{\Figura}[4]{%
    \begin{figure}[H]
        \centering
        \includegraphics[width=#2\textwidth]{#1}
        \caption{#3}
        \label{fig:#4}
    \end{figure}%
}
 desde la raíz del repositorio.
%
%  MOTOR: XeLaTeX (obligatorio). Es el que permite usar la fuente real
%  Times New Roman del sistema. NO compilar con pdflatex.
% =====================================================================

% ---------- Tamaño base exacto: 11 puntos ----------
% Dos precisiones importantes:
%  1) La opción [11pt] de la clase da en realidad 10.95 pt. El paquete
%     fontsize regenera todas las órdenes de tamaño con el valor exacto.
%  2) El "punto" de TeX mide 1/72.27 de pulgada, pero el punto que usan
%     Word y los visores de PDF (el que se verificará) mide 1/72 = 1 bp.
%     Por eso todos los tamaños de este proyecto se declaran en "bp":
%     así el PDF reporta 11.00, 12.00, 14.00 y 16.00 exactos.
\usepackage[fontsize=11bp]{fontsize}

% ---------- Fuentes ----------
\usepackage{amsmath, amssymb}
\usepackage{fontspec}
\defaultfontfeatures{Ligatures=TeX}

% Times New Roman del sistema (macOS y Windows la traen).
% Si no existe, se usa TeX Gyre Termes: clon libre métricamente idéntico,
% así el documento compila igual en cualquier máquina sin cambiar de aspecto.
\IfFontExistsTF{Times New Roman}{
    \setmainfont{Times New Roman}
    \newcommand{\FuenteTexto}{Times New Roman}
}{
    \setmainfont{texgyretermes}[
        Extension      = .otf,
        UprightFont    = *-regular,
        BoldFont       = *-bold,
        ItalicFont     = *-italic,
        BoldItalicFont = *-bolditalic]
    \newcommand{\FuenteTexto}{TeX Gyre Termes (clon métrico de Times New Roman)}
}
% Monoespaciada para el código
\IfFontExistsTF{Courier New}{
    % Sin Scale: así \texttt en línea mide exactamente los 11 bp del texto.
    \setmonofont{Courier New}
}{
    \setmonofont{texgyrecursor}[
        Extension      = .otf,
        UprightFont    = *-regular,
        BoldFont       = *-bold,
        ItalicFont     = *-italic,
        BoldItalicFont = *-bolditalic]
}
% Matemáticas en Times (Termes Math es la versión matemática de Times).
% Se carga por nombre de archivo: en macOS fontconfig no indexa las fuentes
% de TeX Live, así que buscarla por nombre comercial falla.
\usepackage{unicode-math}
\setmathfont{texgyretermes-math.otf}

% ---------- Idioma ----------
% es-lcroman: evita que babel componga los números romanos en versalitas.
% Times New Roman no tiene versalitas reales y XeTeX las simularía
% encogiendo la letra a ~8 pt, rompiendo el requisito de tamaños.
\usepackage[spanish, es-tabla, es-noshorthands, es-lcroman]{babel}
\usepackage{csquotes}

% ---------- Página ----------
\usepackage{geometry}
\geometry{top=2.5cm, bottom=2.5cm, left=2.5cm, right=2.5cm}
\usepackage{microtype}
\linespread{1.15}
\setlength{\parskip}{0.5em}

% ---------- Utilidades ----------
\usepackage{ifthen}
\usepackage{xcolor}
\usepackage{graphicx}
\usepackage{float}
\usepackage{booktabs}
\usepackage{longtable}
\usepackage{multirow}
\usepackage{array}
\usepackage{siunitx}
\usepackage{caption}
\usepackage{subcaption}
\usepackage{enumitem}
\usepackage{tcolorbox}
\usepackage[linesnumbered, ruled, vlined]{algorithm2e}
\usepackage[nottoc, notbib]{tocbibind}   % LOF y LOT aparecen en el índice
\usepackage{fancyhdr}
\usepackage{titlesec}
\usepackage{listings}

% Rutas donde LaTeX buscará imágenes (funciona igual compilando el libro
% completo o un capítulo suelto, porque siempre se compila desde la raíz).
\graphicspath{%
    {img/}%
    {capitulos/cap01/figuras/}%
    {capitulos/cap02/figuras/}%
    {capitulos/cap03/figuras/}%
    {capitulos/cap04/figuras/}%
    {capitulos/cap05/figuras/}%
    {capitulos/cap06/figuras/}%
    {capitulos/cap07/figuras/}%
}

% ---------- Colores institucionales ----------
\definecolor{guindaIPN}{RGB}{110, 28, 52}
\definecolor{azulESCOM}{RGB}{0, 81, 158}
\definecolor{grisCodigo}{RGB}{245, 245, 245}
\definecolor{verdeComent}{RGB}{0, 120, 60}
\definecolor{moradoKey}{RGB}{120, 20, 120}

% =====================================================================
%  JERARQUÍA TIPOGRÁFICA REQUERIDA
%     texto        11 pt
%     capítulo     16 pt negritas
%     sección      14 pt negritas
%     subsección   12 pt negritas
%     subsubsec.   12 pt negritas cursiva
%  Los tamaños se fijan de forma absoluta con \fontsize{pt}{interlínea},
%  no con \Large/\huge, para que sean exactamente los pedidos.
% =====================================================================
\newcommand{\TamCapitulo}{\fontsize{16bp}{19.2bp}\selectfont}
\newcommand{\TamSeccion}{\fontsize{14bp}{16.8bp}\selectfont}
\newcommand{\TamSubseccion}{\fontsize{12bp}{14.4bp}\selectfont}
\newcommand{\TamTexto}{\fontsize{11bp}{13.2bp}\selectfont}

\titleformat{\chapter}[hang]
    {\normalfont\bfseries\TamCapitulo}{\chaptertitlename\ \thechapter.}{0.6em}{}
\titlespacing*{\chapter}{0pt}{0pt}{24pt}

\titleformat{\section}
    {\normalfont\bfseries\TamSeccion}{\thesection}{0.6em}{}
\titlespacing*{\section}{0pt}{16pt}{8pt}

\titleformat{\subsection}
    {\normalfont\bfseries\TamSubseccion}{\thesubsection}{0.6em}{}
\titlespacing*{\subsection}{0pt}{12pt}{6pt}

\titleformat{\subsubsection}
    {\normalfont\bfseries\itshape\TamSubseccion}{\thesubsubsection}{0.6em}{}
\titlespacing*{\subsubsection}{0pt}{10pt}{5pt}

% Encabezados de los índices (son \chapter*) al mismo tamaño de 16 pt
\renewcommand{\contentsname}{Índice de contenido}

% Pies de figura y tabla en el tamaño del texto
\captionsetup{font=normalsize, labelfont=bf, labelsep=period, skip=8pt}

% ---------- Encabezados y pies ----------
\setlength{\headheight}{16pt}
\pagestyle{fancy}
\fancyhf{}
\renewcommand{\headrulewidth}{0.4pt}
\fancyhead[L]{\TamTexto\nouppercase{\leftmark}}
\fancyhead[R]{\TamTexto\Materia}
\fancyfoot[C]{\TamTexto\thepage}
\fancypagestyle{plain}{\fancyhf{}\fancyfoot[C]{\TamTexto\thepage}\renewcommand{\headrulewidth}{0pt}}

% ---------- Código fuente (C, OpenMP, MPI, CUDA) ----------
% El código va en monoespaciada a 10 pt: a 11 pt las líneas largas no caben
% en el ancho de página. El texto redactado sí respeta los 11 pt.
\lstset{
    basicstyle=\ttfamily\fontsize{10bp}{12bp}\selectfont,
    keywordstyle=\color{moradoKey}\bfseries,
    commentstyle=\color{verdeComent}\itshape,
    stringstyle=\color{azulESCOM},
    backgroundcolor=\color{grisCodigo},
    frame=single,
    rulecolor=\color{gray!50},
    numbers=left,
    numberstyle=\fontsize{8bp}{9.6bp}\selectfont\color{gray},
    numbersep=8pt,
    breaklines=true,
    breakatwhitespace=true,
    showstringspaces=false,
    tabsize=4,
    captionpos=b,
    columns=flexible,
}
% Lenguajes derivados
\lstdefinelanguage{CUDA}[ANSI]{C}{
    morekeywords={__global__, __device__, __host__, __shared__, __constant__,
                  __syncthreads, blockIdx, blockDim, threadIdx, gridDim,
                  cudaMalloc, cudaFree, cudaMemcpy, cudaDeviceSynchronize}
}
\lstdefinestyle{C}{language=C, morekeywords={\#pragma, omp, parallel, for, MPI_Init, MPI_Comm_rank, MPI_Comm_size, MPI_Send, MPI_Recv, MPI_Finalize}}
% Nombre en español para los listados de código
\renewcommand{\lstlistingname}{Código}

% ---------- Cajas de nota ----------
\newtcolorbox{nota}[1][Nota]{colback=azulESCOM!5, colframe=azulESCOM!70,
    fonttitle=\bfseries, title=#1, boxrule=0.6pt, arc=2pt}
\newtcolorbox{definicion}[1][Definición]{colback=guindaIPN!5, colframe=guindaIPN!70,
    fonttitle=\bfseries, title=#1, boxrule=0.6pt, arc=2pt}

% ---------- Bibliografía IEEE ----------
\usepackage[backend=biber, style=ieee, sorting=none, maxbibnames=6, minbibnames=3]{biblatex}
\addbibresource{bib/referencias.bib}
\DefineBibliographyStrings{spanish}{references = {Referencias}}

% ---------- Hipervínculos (cargar casi al final) ----------
\usepackage[hidelinks, bookmarksnumbered]{hyperref}
\usepackage{bookmark}

% ---------- Macros de conveniencia ----------
% Bibliografía de la entrega (aparece en el índice de contenido)
\newcommand{\Referencias}{%
    \cleardoublepage
    \printbibliography[heading=bibintoc, title={Referencias}]%
}

% Figura estándar: \Figura{archivo}{ancho}{pie de figura}{etiqueta}
\newcommand{\Figura}[4]{%
    \begin{figure}[H]
        \centering
        \includegraphics[width=#2\textwidth]{#1}
        \caption{#3}
        \label{fig:#4}
    \end{figure}%
}
 desde la raíz del repositorio.
%
%  MOTOR: XeLaTeX (obligatorio). Es el que permite usar la fuente real
%  Times New Roman del sistema. NO compilar con pdflatex.
% =====================================================================

% ---------- Tamaño base exacto: 11 puntos ----------
% Dos precisiones importantes:
%  1) La opción [11pt] de la clase da en realidad 10.95 pt. El paquete
%     fontsize regenera todas las órdenes de tamaño con el valor exacto.
%  2) El "punto" de TeX mide 1/72.27 de pulgada, pero el punto que usan
%     Word y los visores de PDF (el que se verificará) mide 1/72 = 1 bp.
%     Por eso todos los tamaños de este proyecto se declaran en "bp":
%     así el PDF reporta 11.00, 12.00, 14.00 y 16.00 exactos.
\usepackage[fontsize=11bp]{fontsize}

% ---------- Fuentes ----------
\usepackage{amsmath, amssymb}
\usepackage{fontspec}
\defaultfontfeatures{Ligatures=TeX}

% Times New Roman del sistema (macOS y Windows la traen).
% Si no existe, se usa TeX Gyre Termes: clon libre métricamente idéntico,
% así el documento compila igual en cualquier máquina sin cambiar de aspecto.
\IfFontExistsTF{Times New Roman}{
    \setmainfont{Times New Roman}
    \newcommand{\FuenteTexto}{Times New Roman}
}{
    \setmainfont{texgyretermes}[
        Extension      = .otf,
        UprightFont    = *-regular,
        BoldFont       = *-bold,
        ItalicFont     = *-italic,
        BoldItalicFont = *-bolditalic]
    \newcommand{\FuenteTexto}{TeX Gyre Termes (clon métrico de Times New Roman)}
}
% Monoespaciada para el código
\IfFontExistsTF{Courier New}{
    % Sin Scale: así \texttt en línea mide exactamente los 11 bp del texto.
    \setmonofont{Courier New}
}{
    \setmonofont{texgyrecursor}[
        Extension      = .otf,
        UprightFont    = *-regular,
        BoldFont       = *-bold,
        ItalicFont     = *-italic,
        BoldItalicFont = *-bolditalic]
}
% Matemáticas en Times (Termes Math es la versión matemática de Times).
% Se carga por nombre de archivo: en macOS fontconfig no indexa las fuentes
% de TeX Live, así que buscarla por nombre comercial falla.
\usepackage{unicode-math}
\setmathfont{texgyretermes-math.otf}

% ---------- Idioma ----------
% es-lcroman: evita que babel componga los números romanos en versalitas.
% Times New Roman no tiene versalitas reales y XeTeX las simularía
% encogiendo la letra a ~8 pt, rompiendo el requisito de tamaños.
\usepackage[spanish, es-tabla, es-noshorthands, es-lcroman]{babel}
\usepackage{csquotes}

% ---------- Página ----------
\usepackage{geometry}
\geometry{top=2.5cm, bottom=2.5cm, left=2.5cm, right=2.5cm}
\usepackage{microtype}
\linespread{1.15}
\setlength{\parskip}{0.5em}

% ---------- Utilidades ----------
\usepackage{ifthen}
\usepackage{xcolor}
\usepackage{graphicx}
\usepackage{float}
\usepackage{booktabs}
\usepackage{longtable}
\usepackage{multirow}
\usepackage{array}
\usepackage{siunitx}
\usepackage{caption}
\usepackage{subcaption}
\usepackage{enumitem}
\usepackage{tcolorbox}
\usepackage[linesnumbered, ruled, vlined]{algorithm2e}
\usepackage[nottoc, notbib]{tocbibind}   % LOF y LOT aparecen en el índice
\usepackage{fancyhdr}
\usepackage{titlesec}
\usepackage{listings}

% Rutas donde LaTeX buscará imágenes (funciona igual compilando el libro
% completo o un capítulo suelto, porque siempre se compila desde la raíz).
\graphicspath{%
    {img/}%
    {capitulos/cap01/figuras/}%
    {capitulos/cap02/figuras/}%
    {capitulos/cap03/figuras/}%
    {capitulos/cap04/figuras/}%
    {capitulos/cap05/figuras/}%
    {capitulos/cap06/figuras/}%
    {capitulos/cap07/figuras/}%
}

% ---------- Colores institucionales ----------
\definecolor{guindaIPN}{RGB}{110, 28, 52}
\definecolor{azulESCOM}{RGB}{0, 81, 158}
\definecolor{grisCodigo}{RGB}{245, 245, 245}
\definecolor{verdeComent}{RGB}{0, 120, 60}
\definecolor{moradoKey}{RGB}{120, 20, 120}

% =====================================================================
%  JERARQUÍA TIPOGRÁFICA REQUERIDA
%     texto        11 pt
%     capítulo     16 pt negritas
%     sección      14 pt negritas
%     subsección   12 pt negritas
%     subsubsec.   12 pt negritas cursiva
%  Los tamaños se fijan de forma absoluta con \fontsize{pt}{interlínea},
%  no con \Large/\huge, para que sean exactamente los pedidos.
% =====================================================================
\newcommand{\TamCapitulo}{\fontsize{16bp}{19.2bp}\selectfont}
\newcommand{\TamSeccion}{\fontsize{14bp}{16.8bp}\selectfont}
\newcommand{\TamSubseccion}{\fontsize{12bp}{14.4bp}\selectfont}
\newcommand{\TamTexto}{\fontsize{11bp}{13.2bp}\selectfont}

\titleformat{\chapter}[hang]
    {\normalfont\bfseries\TamCapitulo}{\chaptertitlename\ \thechapter.}{0.6em}{}
\titlespacing*{\chapter}{0pt}{0pt}{24pt}

\titleformat{\section}
    {\normalfont\bfseries\TamSeccion}{\thesection}{0.6em}{}
\titlespacing*{\section}{0pt}{16pt}{8pt}

\titleformat{\subsection}
    {\normalfont\bfseries\TamSubseccion}{\thesubsection}{0.6em}{}
\titlespacing*{\subsection}{0pt}{12pt}{6pt}

\titleformat{\subsubsection}
    {\normalfont\bfseries\itshape\TamSubseccion}{\thesubsubsection}{0.6em}{}
\titlespacing*{\subsubsection}{0pt}{10pt}{5pt}

% Encabezados de los índices (son \chapter*) al mismo tamaño de 16 pt
\renewcommand{\contentsname}{Índice de contenido}

% Pies de figura y tabla en el tamaño del texto
\captionsetup{font=normalsize, labelfont=bf, labelsep=period, skip=8pt}

% ---------- Encabezados y pies ----------
\setlength{\headheight}{16pt}
\pagestyle{fancy}
\fancyhf{}
\renewcommand{\headrulewidth}{0.4pt}
\fancyhead[L]{\TamTexto\nouppercase{\leftmark}}
\fancyhead[R]{\TamTexto\Materia}
\fancyfoot[C]{\TamTexto\thepage}
\fancypagestyle{plain}{\fancyhf{}\fancyfoot[C]{\TamTexto\thepage}\renewcommand{\headrulewidth}{0pt}}

% ---------- Código fuente (C, OpenMP, MPI, CUDA) ----------
% El código va en monoespaciada a 10 pt: a 11 pt las líneas largas no caben
% en el ancho de página. El texto redactado sí respeta los 11 pt.
\lstset{
    basicstyle=\ttfamily\fontsize{10bp}{12bp}\selectfont,
    keywordstyle=\color{moradoKey}\bfseries,
    commentstyle=\color{verdeComent}\itshape,
    stringstyle=\color{azulESCOM},
    backgroundcolor=\color{grisCodigo},
    frame=single,
    rulecolor=\color{gray!50},
    numbers=left,
    numberstyle=\fontsize{8bp}{9.6bp}\selectfont\color{gray},
    numbersep=8pt,
    breaklines=true,
    breakatwhitespace=true,
    showstringspaces=false,
    tabsize=4,
    captionpos=b,
    columns=flexible,
}
% Lenguajes derivados
\lstdefinelanguage{CUDA}[ANSI]{C}{
    morekeywords={__global__, __device__, __host__, __shared__, __constant__,
                  __syncthreads, blockIdx, blockDim, threadIdx, gridDim,
                  cudaMalloc, cudaFree, cudaMemcpy, cudaDeviceSynchronize}
}
\lstdefinestyle{C}{language=C, morekeywords={\#pragma, omp, parallel, for, MPI_Init, MPI_Comm_rank, MPI_Comm_size, MPI_Send, MPI_Recv, MPI_Finalize}}
% Nombre en español para los listados de código
\renewcommand{\lstlistingname}{Código}

% ---------- Cajas de nota ----------
\newtcolorbox{nota}[1][Nota]{colback=azulESCOM!5, colframe=azulESCOM!70,
    fonttitle=\bfseries, title=#1, boxrule=0.6pt, arc=2pt}
\newtcolorbox{definicion}[1][Definición]{colback=guindaIPN!5, colframe=guindaIPN!70,
    fonttitle=\bfseries, title=#1, boxrule=0.6pt, arc=2pt}

% ---------- Bibliografía IEEE ----------
\usepackage[backend=biber, style=ieee, sorting=none, maxbibnames=6, minbibnames=3]{biblatex}
\addbibresource{bib/referencias.bib}
\DefineBibliographyStrings{spanish}{references = {Referencias}}

% ---------- Hipervínculos (cargar casi al final) ----------
\usepackage[hidelinks, bookmarksnumbered]{hyperref}
\usepackage{bookmark}

% ---------- Macros de conveniencia ----------
% Bibliografía de la entrega (aparece en el índice de contenido)
\newcommand{\Referencias}{%
    \cleardoublepage
    \printbibliography[heading=bibintoc, title={Referencias}]%
}

% Figura estándar: \Figura{archivo}{ancho}{pie de figura}{etiqueta}
\newcommand{\Figura}[4]{%
    \begin{figure}[H]
        \centering
        \includegraphics[width=#2\textwidth]{#1}
        \caption{#3}
        \label{fig:#4}
    \end{figure}%
}
